% ============================================================================
%  Chapter 5. Discussion
%  \input file for a 12pt double-spaced `report' document.
%  Preamble must provide: booktabs, float (for [H]), amsmath, natbib.
%  Bibliography: accord/thesis/refs.bib
%  Numeric provenance: every quantity is from results/accord/VERIFIED_FACTS.json.
%  Compute figures come from results/accord/control_cost_log.json, the source
%  named for that item alone. No open markers remain in this chapter.
%
%  Figures: fig:forest (evaluation survived, second family), fig:scaling (the null is a
%  measurement). Sources for each are in figures/manifest.json and the
%  generating scripts beside it. The actionability timeline and the exchange-rate
%  figure were removed; the sentence and paragraph that carried their claims stayed
%  in prose. The threshold-surface figure for the
%  selection-channel anomaly moved with the rest of Section 5.3's detail to
%  the supplementary file; see the pointer at the end of that section.
% ============================================================================

\chapter{Discussion}
\label{ch:discussion}

Chapter~\ref{ch:results} reported what each stage of the pipeline measured, and
what the control attached to that stage measured afterwards. This chapter argues
about the difference. Three claims are defended: the controls changed the reading
at every stage and reversed it at several; the training and evaluation harness
resolves real differences at the sample sizes it was given, and the
twelve-positive contradiction evaluation's own sensitivity, fixed by its class
counts rather than by the harness, is a separate and narrower question the
harness result does not by itself answer; and one result here is unexplained
and should stay unexplained
until somebody tests a mechanism rather than proposing one.

\section{What Each Control Changed}
\label{sec:disc-spine}

Table~\ref{tab:disc-spine} is the summary this section argues through.

The pipeline has four stages and a reporting layer. Each produced a number a
competent paper could have reported, and no standard metric at that stage was
sensitive to the defect sitting behind the number; each was then subjected to
a control, cheap against what it checked, that was. The controls are not new: negative
controls come from observational epidemiology \citep{lipsitch2010negative}, with the positive
direction adapted here rather than taken from that source,
randomisation sanity checks from the interpretability literature
\citep{adebayo2018sanity}, control tasks and selectivity from probing
\citep{hewitt2019control}, and auditing a pipeline stage by stage has been
demonstrated for fairness in tabular preprocessing
\citep{biswas2021fairpreprocessing}. What is new is that they were run here, and
what they returned.

\begin{table}[H]
\centering
\caption[Each pipeline stage before and after control]{Each pipeline stage, the reading available before its control, and the
reading after. ``Reversed'' means the direction of the conclusion changed.
``Narrowed'' means it survived in a smaller form. ``Survived'' means the control
fired and the original reading held.}
\label{tab:disc-spine}
\small
\begin{tabular}{p{0.18\linewidth}p{0.25\linewidth}p{0.25\linewidth}p{0.14\linewidth}}
\toprule
\textbf{Stage} & \textbf{Before control} & \textbf{After control}
& \textbf{Verdict} \\
\midrule
Candidate mining & Pool base rate $0.0401$ read as a property of the mining rules & $0.0401$ is an upper bound on annotation validity, not an estimate of pool quality & Narrowed \\
\addlinespace
Weak supervision & A high-recall labeller usable in place of expert annotation & $261$ of $299$ gold pairs labelled positive against $12$ true positives; a prior, not a decision & Reversed \\
\addlinespace
Architecture & Three named components credited with the reported performance & Three different failures, not one: the authority scalars are falsified, sitting at initialisation across seeds; the rule-type head is confounded, its metric vacated by a target one linear layer from an input feature that equals it; the harm head is vacuous, with no term in the objective to train it. No randomisation control was run on any of the three & Reversed \\
\addlinespace
Evaluation & AUROC near chance, plausibly an artifact of a small and defective sample & At chance on the defect-free stratum, with an instrument shown to resolve $4.53$ points where an effect existed & Survived \\
\addlinespace
Reporting & Drift slopes read directionally from stored gap magnitudes & Same magnitudes, reproduced to a maximum discrepancy of $0.0$; roughly half the slope signs invert & Reversed \\
\bottomrule
\end{tabular}
\end{table}

\subsection{Candidate mining: narrowed}

The miner produced a pool whose base rate against pharmacist annotation is $0.0401$, $12$ true
contradictions in the $299$-pair full annotated set; on the $253$-pair decontaminated core the same
$12$ positives sit against $241$ negatives for a base rate of $0.0474$, and ACCORD thresholded at
$0.5$ on that core returns $11$ true positives against $223$ false positives, a precision of
$0.047$ that is the detector's confusion matrix rather than a further reading of pool quality.
Neither figure is a property of the annotation exercise on its own: for action-conflict candidates
the trigger quantified over sets of medication mentions while the exemplar selector indexed
positionally, and the audit found the two did not reliably agree. A reconstruction of the exemplar,
at $0.8934$ to $0.9508$ fidelity, resolves $137$ of the $165$ Rule 1 and Rule 2 candidates and finds
$65$ wrong, a rate of $0.4745$; a true key join covers only $23$ of those rows and finds all $23$
correct. The two cannot both describe the same rows, so they are reported as mutually exclusive
rather than as a range \citep{jacobs2021measurement}, and the direction of the resulting bound on
pool precision is unknown: the pharmacist judged the material shown, not corrected material, so a
higher and a lower true precision both remain consistent with the record. One stratum escapes the
defect entirely, Rule 3 and Rule 4, correct by construction, holding $134$ pairs and $6$ of the
$12$ confirmed contradictions, corroborating $0.0401$ on ground the exemplar defect cannot reach.

\subsection{Weak supervision: a prior confirmed locally, a scaling claim now withdrawn}

The labeller looked usable on the property it was checked against, recall. The control is the one
Feinstein and Cicchetti described for a different instrument thirty-six years ago
\citep{feinstein1990paradoxes,cicchetti1990resolving}: compare the marginal against a
human-established base rate before believing anything about agreement. Against the $299$ annotated
pairs the deployed local labeller assigned the positive class to $261$; the pharmacist assigned it
to $12$. Saying yes to $87\%$ of a pool whose true positive rate is $4\%$ is not a decision, it is a
prior with a text prompt attached, and three further local eight-billion-parameter models reproduce
the same miscalibration at over-calling ratios from $15.8196$ to $23.5959$
(Chapter~\ref{ch:results}), so this is not one model misbehaving. Correcting a marginal this biased
needs gold labels on the same distribution \citep{forman2008quantifying,egami2023surrogates}, which
a workflow built to avoid annotation cannot supply, so the student inherits the marginal rather
than the task \citep{ratner2016dataprogramming,ratner2017snorkel,zhang2022wssurvey}. What this
finding cannot now say is how that prior scales with capability, vendor or serving channel: the
comparisons built to answer that sent sentence pairs to hosted services, and every number they
produced has been withdrawn with the comparisons themselves (Section~\ref{sec:m-ethics}); the
question is indeterminate rather than resolved in either direction. The surviving claim is
narrower than the one this section once made: not that language models cannot reason about
clinical text, but that this project's weak labeller's marginal was a property nobody checked
before training on it \citep{pangakis2023validation}.

\subsection{Architecture: three components, three different failures}

The architectural checks did the most damage, and the word for what they did differs by component,
because no randomisation control was run on any of the three clinical components; the
randomised-substitute arms in this thesis belong to a second, provenance-conditioned system and do
not transfer here. Only the authority scalars are falsified in the strict sense: the parameter is
initialised at the value later reported as discovered, sits behind a normalisation that removes its
gradient, and agrees across seeds, so the claim that training found it is refuted. The rule-type
head is confounded rather than falsified, its auxiliary target readable from the input features one
linear layer away, the textbook symptom of leakage rather than learning \citep{kaufman2012leakage},
which vacates the metric rather than refuting the component. The harm head is neither, because no
term in the objective trains it, so no mechanism ever asserted anything a control could refute, and
its severity claim is withdrawn as vacuous. A fourth check, targeting the pretraining collapse
rather than the architecture, tested whether anisotropy rather than clinical homogeneity explains a
high similarity score: unrelated non-clinical sentences return a mean pairwise similarity of
$0.9113$ against $0.9176$ on matched clinical pairs, a Welch test rejecting exact equality
($t = 4.0001$, Cohen's $d = 0.2808$), which undercuts the homogeneity account without establishing
identity. A component credited by a system's own authors and never given a substitute that should
not work, the control-task logic \citet{hewitt2019control} developed and this thesis applies rather
than invents, has simply not been tested.

\subsection{Evaluation: survived}

The detector's ranking is at chance and stays there under every control applied to it, which is why
this thesis has a finding rather than only a list of retractions. The decontaminated ACCORD
ensemble reaches AUROC $0.5145$ ($p = 0.8668$) on the $253$ non-ambiguous pairs against $0.4595$
for the contaminated released run; the leaked rows carried the weak labeller's labels rather than
the pharmacist's, so the leak had no channel to inflate agreement with gold, making the contaminated
figure a conservative rather than an inflated comparator, and the $0.0550$ gap carries no
significance claim since no paired comparison exists in the released record. A separate,
defect-free control model, a plain BioMedBERT with none of ACCORD's authority parameters, rule
head or harm head, reaches a three-seed mean AUROC of $0.4967 \pm 0.0350$ on all $299$ pairs
(ensemble $0.4692$) and $0.5295 \pm 0.0158$ on the $134$-pair Rule 3 and Rule 4 stratum (ensemble
$0.5469$); ACCORD's own contaminated run, rescored there, reads $0.4602$ and $0.5078$, so the
control ensemble sits $0.0090$ above it at $p = 0.9300$ on the full set and $0.0391$ above it at
$p = 0.4306$ on the stratum, every interval straddling chance (Figure~\ref{fig:forest}, second
family). The stratum check has almost no power, since $6$ positives put its minimum detectable
AUROC at $0.8058$ and only the contaminated ACCORD score exists there, so the objection that real
signal was attenuated by wrong-evidence labels is not answered directly. Two baselines finish the
picture: a majority-class predictor gives AUROC $0.5$ and average precision $0.0401$, the base
rate, and a TF-IDF logistic regression gives AUROC $0.5761$, average precision $0.1571$, placing
$2$ true positives in its top $15$ where the neural detector places $0$; neither model is useful,
but no reading survives in which the custom architecture contributed something a bag-of-words
baseline could not.

\subsection{What survived}

Two things survived their controls, and a section in which nothing survives is not a finding. The
chance-level ranking survived decontamination, seed variation, defect-free stratification, a
DeLong comparison against an ensemble, and two baselines. The over-calling survived three changes
of local model family; whether it also survives a change of scale or serving vendor is now
indeterminate rather than confirmed, because the comparisons built to test that sent sentence pairs
to a hosted channel and every number they produced has been withdrawn (Section~\ref{sec:m-ethics}).
The direction of both surviving conclusions held through every correction applied to them, each of
which moved a headline number toward the negative reading rather than away from it. What did not
survive, and what the withdrawal narrows further still, is any claim about mechanism or about
scale: this is a measurement of a pipeline, not an account of why clinical text is hard or of how
the failure would change at a different model size.

\section{What the Positive Control Licenses, and What It Does Not}
\label{sec:disc-null}

A negative result in a small clinical sample invites one objection above all
others: the experiment was too small to see anything, so it saw nothing. That
objection is correct about most such results, and would be correct about this one
if no positive control existed. One does. The device adapts the one
\citet{lipsitch2010negative} formalise for observational research: run the
analysis on a case where the answer is known in advance, and let the result
speak about the instrument rather than about the substantive question.

The instrument is a comparison of a domain-pretrained encoder against an
identically sized encoder trained from scratch, evaluated by a convergence sweep
of $48$ runs. On a public four-class news classification task it resolves a
pretraining benefit of $4.53$ points at $24$ epochs: $0.9379$ for the pretrained
arm at SD $0.0029$ against $0.8925$ for the scratch arm at SD $0.0004$. The bands
are disjoint by a wide margin. At $3$ epochs the gap is $5.07$ points, so the
effect is present early and does not depend on the budget chosen.

The budget is what closes the objection. Eight times the training budget bought
the 100M-parameter scratch arm $0.41$ points, from $0.8885$ at $3$ epochs to $0.8925$
at $24$; the pretrained arm moved $-0.13$ points over the same range. Both arms
are converged. A gap that persists at convergence is not a gap in optimisation,
and an instrument reporting a $4.53$-point difference at standard deviations
at or below $0.0036$ is not an instrument that cannot see effects.

Parameters do not close it either. Table~\ref{tab:scaling} reports the from-scratch
arms at three sizes against the pretrained band on the public task. The
from-scratch curve rises with size and does not approach the band, so the public
gap is a pretraining effect rather than a capacity effect, and the mortality
control returns almost nothing (Figure~\ref{fig:forest}, last row).


The decision rule was committed before the sweep ran, at repository commit
\texttt{b048829}, and the outcome it adjudicated
is the one that keeps the claim rather than the one that discards it. That
ordering is the whole value of pre-registration, and it could have gone the other
way: arms that had not converged would have retired the positive control and with
it the licence to call the mortality-control result a measurement. On the
mortality control task, in-hospital mortality predicted from a patient's first
$48$ hours of notes on $6{,}776$ held-out admissions holding $570$ positives, the
same instrument, harness and protocol return $+0.0031$, with the pretrained arm
at $0.7693$ (SD $0.0102$) against $0.7662$ (SD $0.0073$), each at three seeds.
No significance test for that difference is registered, so no $p$ value is
quoted; the two standard deviations alone make the gap look like noise around a
shared value rather than a resolved difference, and the sign is not robust:
restricting the from-scratch arm to the two runs whose data lineage is recorded
gives $0.7703$ and a gap of $-0.0011$. Both readings are nulls and the argument
below turns on neither. Same instrument, same protocol, two tasks, one
showing $4.53$ points and the other showing nothing.

The general result is not this thesis's. Pretraining loss is a poor proxy for
downstream performance, and that is settled at far larger scale than $48$ runs
\citep{lourie2025scaling,feng2026tessera}. The
sweep here is a positive control for the mortality-task measurement, and writing
it up as a contribution would misrepresent what it is for.

What this control licenses is narrower than it can sound: that the harness
resolves a real difference at the sample sizes it was given, not the
sensitivity of the twelve-positive contradiction evaluation, a different
instrument with its own minimum detectable AUROC, $0.7355$ at $12$ positives
against $241$ negatives ($0.7350$ on the $299$-pair set, $0.8058$ on the
$6$-positive defect-free stratum), fixed by that evaluation's class counts and
corroborated arithmetically by simulation in Chapter~\ref{ch:results}. This work
therefore cannot establish that cross-note medication contradiction detection
is unsolvable; it establishes that this pipeline does not solve it, and that an
effect at or above $0.7355$ would have been visible had one been present. Only
the second statement is made.

A bootstrap standard deviation of $0.0870$ on the contaminated released run's
$G_{299}$ AUROC is consistent
with a mildly informative model and with a mildly inverted one. The three-seed
ensemble's own $95\%$ bootstrap interval on the $253$-pair set,
$[0.3210, 0.7034]$, is wide enough that its upper bound of $0.7034$ sits only
$0.0321$ below the minimum detectable AUROC of $0.7355$, so interval width on its own does not rule out a
modestly useful model. Three of the four estimators compared here, the single
model, the three-seed ensemble and the TF-IDF baseline, could in principle have
found separation and did not; the fourth, a majority-class predictor, returns
AUROC $0.5$ by construction and was never in a position to look, so it is a
reference point rather than a fourth line of evidence. What the three estimators
that could fail and did not agree on is the evidence a sample this size can
produce, not proof that no separation exists.

\section{An Anomaly, Reported as an Anomaly}
\label{sec:disc-anomaly}

One result here is not explained, and this section reports it without
explaining it. It is not counted among Chapter~\ref{ch:intro}'s
contributions, and that billing is deliberate rather than an oversight: a
contribution earns its place by explaining something, and this pattern
explains nothing about itself. The population audited is simulated, not the
clinical pipeline; it is retained because the same experiments also anchor
the sign-provenance rule of Chapter~\ref{ch:results}, applied afterward to
the clinical numbers, and because a pattern that does not fit is better
reported than dropped.

On two simulated populations, with baseline disparity held fixed, upstream
selection channels were swept and the pool-level demographic parity gap
audited at each step. The gap never changes sign within a sweep, but the
direction it drifts in does: two channels swept on the same population give
fitted drift slopes of opposite sign, and on the population where the sweep
was extended across operating points as well as channel strength, the
audited gap moves with both, so a single audited number describes one cell of
a surface rather than the system that produced it. A further, controlled
sweep of the same construct across baselines, channel strengths and
replicates does not reproduce even that channel dependence: it inverts the
mapping between channel and drift direction, so the direction a pool audit
drifts in is not a stable function of the channel that produced it, at least
across the region examined.

Five mechanism claims were proposed during this project to account for this
pattern; each was tested against a parameter sweep and each failed, and a
sixth is not offered. No subgroup harm is demonstrated, none is demonstrated
anywhere in this work, and no general claim is made that disaggregated
fairness evaluation is uninterpretable under selection; that argument belongs
to others and predates this work
\citep{pfohl2025disaggregated,fogliato2020fairness}, as does the observation
that the design choices inside a disaggregated evaluation are consequential
enough to need documenting \citep{barocas2021disaggregated}. The full sweep,
the operating-point surface, and the failed mechanism accounts are given in
Appendix~\ref{app:anomaly}.

\section{Measurement Provenance}
\label{sec:disc-provenance}

Every disparity gap in this project was stored as an absolute value, and most
ratio quantities as a minimum over a maximum, with no per-group values persisted.
Each decision is individually reasonable: a magnitude is the right object for a
pass or fail check against a tolerance. Together they destroy the information
every directional sentence in the analysis depended on.

The retrofit is the useful part. Signed storage was added alongside the existing
fields rather than replacing them, and the recomputation reproduced every prior
value with a maximum discrepancy of $0.0$; the reproduction check passes. Nothing
previously reported changed by any amount. What changed was the sign of roughly
half the drift slopes, and one credibility call: a slope whose stored interval
crossed zero has a signed interval that does not. The flip count reported here, $7$ of $14$
compared metric-slopes, is the figure carried into the verified register
(Appendix~\ref{app:corrections}, C57). The field count behind the retrofit is not registered and
is not quoted.

The rule is one sentence. A metric stored as a magnitude licenses no directional
claim, and no downstream care recovers what the storage discarded. It is not
novel and should not be dressed as novel; it is a specific case of a general
point about measurement provenance and about the debt that accumulates when an
intermediate representation is chosen for one consumer and read by another
\citep{sculley2015debt}. What this work adds is a demonstrated true positive.
Applied to this project's own artifacts, the check invalidated four directional
readings, three of which were themselves corrections of earlier errors. A rule
that catches a correcting analysis committing the error it is correcting has some
claim on a practitioner's attention. Two qualifications: the check is cheap only
because these sweeps were cheap, and where per-group values were never written
the retrofit is a re-run rather than a recomputation; and the demonstration is on
one project's artifacts, so no claim is made about how often the defect occurs
elsewhere.

\subsection{What practice does with the upstream check}

The survey speaks to whether the check that would have caught the upstream defect
is run at all. Auditing a pipeline one stage at a time is not an exotic ask:
\citet{biswas2021fairpreprocessing} do it for tabular preprocessing operators,
and show that the local behaviour of each stage composes into the behaviour of
the pipeline. On a $75$-paper corpus, $4$ papers audit the selection step under
the strict reading and $14$ under the permissive one, so $61$ report it under
neither. What those $61$ do instead is not recorded by this coding, which asks
only whether the check appears; a separate coded field, how each paper handles
null attribute values, is reported over the same $75$ papers in
Table~\ref{tab:survey}, and the two fields are different variables on the same
corpus rather than a partition of it.

Both readings appear because the coding cannot separate them cleanly, and
reporting only the one that suits the argument would be the error this thesis is
about. The corpus was deliberately enriched for fairness-forward work, and the
enrichment runs against the finding rather than for it: a corpus over-sampling
the papers most likely to run the check still returns $4$ of $75$ strictly. This
is not a prevalence estimate for clinical machine learning. Coverage is
open-access and preprint venues rather than the literature as a whole, and
$459$ screened-in papers were never retrieved in full text.

\section{What the Controls Cost}

Table~\ref{tab:disc-cost} itemises it.
\label{sec:disc-cost}

Nobody reports this, so it is reported here, with the limit stated first.
\textbf{These figures are compute only.} Person-hours were not logged for most
of the project and are not reconstructed, because a reconstructed estimate
would be an invented number in a thesis about invented numbers.

\begin{table}[H]
\centering
\caption[Recorded compute cost of controls]{Recorded compute for the controls. Person-hours are excluded and were
not logged. The cost log records a $12$-run, $34.9$-minute public-task episode
but omits the $48$-run, $555$-minute convergence sweep the argument in
Section~\ref{sec:disc-null} rests on. The two are separate run sets, not one
nested in the other: nine of the twelve cells differ from the convergence
sweep's shortest-budget cells and all twelve recorded runtimes differ, so both
are counted here.}
\label{tab:disc-cost}
\begin{tabular}{lrr}
\toprule
\textbf{Control episode} & \textbf{Runs} & \textbf{GPU-minutes} \\
\midrule
Public-task convergence sweep             & 48 & 555.0 \\
Public task, single budget across scales  & 12 & 34.9 \\
Clinical task, pretrained against scratch & 17 & 612.3 \\
Signed re-run, selection experiment       & 1  & 0.7 \\
Signed re-run, PUMS replication           & 1  & 13.5 \\
\midrule
Total, reconciled                         & 79 & 1216.4 \\
\bottomrule
\end{tabular}
\end{table}

The control that retracted the directional core of three corrections cost
$0.7 + 13.5 = 14.2$ GPU-minutes across both experiments. The convergence sweep that licenses the
whole argument of Section~\ref{sec:disc-null} cost $555$ minutes on its own. Together with the
clinical-task comparison the two larger sweeps cost $555 + 612.3 = 1167.3$
GPU-minutes, and the reconciled total is $555 + 34.9 + 612.3 + 14.2 = 1216.4$
GPU-minutes and $79$ runs in total, roughly twenty GPU-hours. The claim is not that
controls are always cheap: re-annotation is not, a second annotator is not, and
the one control here that cannot be recomputed from disk requires re-running the
experiment that produced it. The claim is narrower. In this project the controls
that reversed conclusions cost less compute than the experiments whose
conclusions they reversed, and cost is not why they went unrun.

No comparison against prose checklists is offered. An earlier version of this
work compared executable gates against checklist-style review, and that
comparison was retracted: the gate recorded no failure in any run, so its catch
rate is undefined, while the checklist arm computed the relevant defect, wrote it
to disk, and was misread by a human. Executability was not the binding
constraint, and no claim is made that executable gates outperform prose
checklists. The nearest thing to evidence on that question comes from journal
peer review rather than from here, where introducing an automated statistical
consistency check is associated with a steep decline in reporting
inconsistencies \citep{nuijten2023statcheck}, measured over a corpus this
project cannot approach.

\section{Limitations}
\label{sec:disc-limits}

\paragraph{RQ4's answer, now that the hosted channel is ruled out.} One dependency in this thesis
ran through four chapters and is stated here as a settled finding rather than a hypothetical,
because the question it depended on has been answered rather than left open. Section~\ref{sec:m-ethics}
records that every hosted-service transmission this project made, on the Anthropic API and on AWS
Bedrock alike, has been withdrawn: PhysioNet's guidance of 24 September 2025 prohibits third-party
retention without a vendor qualifier, and every hosted route used here transmitted credentialed
sentence pairs off the analysis machine, so every number those runs produced has been removed from
the set a reported figure may resolve against. Those rows, not the local eight-billion-parameter
models, were the evidence that the twelve-positive contradiction evaluation could resolve an effect
at all, so Section~\ref{sec:r-frontier}'s demonstration reverts to the calculation it was built to
replace, and RQ4's answer moves from ``the instrument was sensitive enough to see an effect of that
size'' to ``the minimum detectable effect says an effect of that size should have been visible, and
nothing that remains in the register confirms that one was.'' No vendor-specific fallback survives
this: an earlier reading of this dependency distinguished Amazon Bedrock, on the retention test
PhysioNet's guidance sets out, from the Anthropic API, but the withdrawal recorded in
Section~\ref{sec:m-ethics} applies uniformly to every hosted route used here, so that distinction no
longer has anything left to save. What still stands is the pretraining positive control, which
shows the same measurement apparatus resolves a real difference on a public benchmark unconnected
to any hosted transmission, and the rule-head transfer collapse, neither of which touches a hosted
API. That is a weaker position than an earlier version of this chapter argued from, and it is the
position this thesis now holds.

The limitations below are stated without hedging. Several are severe enough that
a reader could reasonably decline to draw a conclusion from part of this work,
and where that is so it is said.

\paragraph{Twelve positives.} Every ranking measurement rests on $12$ positive
cases, and the defect-free stratum on $6$. The minimum detectable AUROC at $80\%$
power is $0.7355$ on the primary set and $0.8058$ on the stratum. The intervals
are wide and the ordering between the two systems this pipeline produced is
unidentified, their Matthews correlations being $-0.0436$ and $-0.0070$, both
within $0.05$ of chance. The limit is not that $12$ positives cannot separate
anything: two of the local weak-labeller replicates from
Section~\ref{sec:m-control-rq2}, \texttt{granite3.3:8b} zero-shot and
\texttt{llama3:8b} few-shot, carry Matthews correlation intervals against these
same $253$ rows that exclude zero, at $0.1326$ and $0.0540$. That is a claim
about those two labeller configurations on the agreement-with-gold question, not
about the contradiction detector, and the qualification belongs in the same
breath as the claim rather than later: $12$ positives can separate a system from
chance when the effect is large enough, so the null the detector and its control
return is a statement about those two architectures, which sit close to chance
and to each other, rather than a statement that no instrument at this sample
size could ever resolve anything. No ablation carries a significance claim.

\paragraph{Deployment risk, and automation bias.} A system with recall
$0.7778$ and precision $0.0424$ generates a large number of false alarms per
true contradiction: the decontaminated detector's first seed alone raises
$216$ false alarms against $9$ true contradictions. Alert burdens at that
ratio are associated with reduced acceptance
\citep{vandersijs2006overriding,ancker2017alertfatigue}, so this system should
not be deployed in any form; its failure mode is worse than doing nothing, since it consumes
clinician attention while providing no information. A second, more specific
hazard sits inside the architecture: the harm head of
Section~\ref{sec:m-detector} received no loss term and its outputs span
$0.3977$ to $0.6096$, an untrained sigmoid's range, yet an earlier version of
this work presented that output as a severity score, a plausible danger
score derived from random initialisation, the specific and most serious
instance of automation bias this project produced. The diagnostic that
catches it, printing the empirical range of every emitted score, costs one
line.

\paragraph{One institution, one corpus, one era.} MIMIC-III covers admissions to
a single United States academic medical centre \citep{johnson2016mimic}. The
primary cohort is a billing-code phenotype, over-including admissions where the
condition is a comorbidity and under-including those where it was present and
uncoded. Documentation conventions, note-type taxonomies and medication-list
formatting are institution- and era-specific, and the drug dictionary, class
dictionary, dose threshold and danger judgements all encode that context. The
transfer cohorts move across index diagnosis within the same institution, so
nothing speaks to cross-institutional generalisation
\citep{zech2018generalization}. MIMIC-III is a credentialed PhysioNet resource;
all analysis reports aggregates only, and no clinical note text and no
identifiers appear in this document or in any released artifact.

\paragraph{The trigger fires late, and the record cannot say how late.} The
system compares a nursing note against a discharge summary, so it produces
nothing until the discharge summary exists, and that document closes the
admission rather than preceding it, finalised at or after discharge rather than
before it. The data model compounds the problem: nursing notes carry a
charting timestamp and discharge summaries carry only a charting date, so the
summary cannot be placed in time relative to the discharge itself. That same
field absence is what invalidated an earlier prefilter result in this project,
retracted in Appendix~\ref{app:corrections}. A contradiction detected at that
point is a documentation finding rather than an intervention, and a deployment
claim would have to move the trigger earlier, to the discharge order rather than
to the finished summary. Nothing here measures whether the task is tractable at
that earlier point, and nothing here should be read as a claim that it is.

\paragraph{The exemplar defect bounds annotation validity, and re-annotation was
not performed.} This is the limitation most likely to be read as smaller than it
is. The rebuild that corrects the selector exists; the re-annotation that would
convert the bound of Section~\ref{sec:disc-spine} into an estimate was not run.
So $0.0401$ cannot be cited as a property of the mining rules by anyone, this
thesis included, and a reader who wants that quantity should treat it as
unmeasured rather than as measured with error. What the defect does not touch is
the central evaluation result, which reproduces on the $134$-pair stratum. The
defect-free stratum, correct by construction, offers a check that does not
depend on that estimate: it holds $6$ of the $12$ confirmed contradictions in
$134$ pairs, a rate of $0.0448$, corroborating $0.0401$ on a stratum where the
exemplar defect cannot fire. The weak-supervision finding needs no base-rate
estimate at all: $261$ of the $299$ gold pairs were labelled positive against
$12$ confirmed, a marginal comparison with no denominator to estimate.

\paragraph{The survey is a prevalence measurement on an enriched corpus, and it
is limited on several axes at once.} It is the weakest of the empirical
strands in this thesis, reported at that weight rather than alongside the
ablation or the pre-registered convergence control. The $75$-paper corpus was
assembled to over-sample fairness-forward work, coverage is open-access
sources only, $459$ screened-in papers were never retrieved in full text, and
one coder codes the full set; Chapter~\ref{ch:results} gives the full
enrichment argument behind $4/75$ and $14/75$. A blind second coder recoded a
subsample with field definitions only and the hypothesis withheld, and
agreement on the permissive field is high. On the 19 double-coded papers, the
registered permissive
agreement is $0.9474$ with kappa $0.855$. On the strict cell, the one that decides the
headline count, the reliability coefficient is
\textbf{undefined}: both coders returned the same code on all 19 overlap
papers, which drives a coder's marginal on the category toward degenerate and
the denominator of a chance-corrected coefficient toward zero. The full
contingency table behind the undefinedness is not registered and is not
claimed here beyond that. It is reported as undefined rather than as the
perfect agreement the raw agreement rate superficially permits, which is the
same class of error as reporting a coefficient guaranteed by construction.
The survey coding table of Appendix~\ref{app:survey} reproduces the full 75-paper coding, so a
reader who wants to check the count rather than take it on trust can.

\paragraph{One annotator, and the reliability figure this instrument cannot
produce.} Every gold label came from a single board-certified
pharmacist in two passes separated by a washout. The washout controls
within-session drift; it does not establish inter-rater reliability, because
there is one rater. A previously reported perfect agreement between two
``independent judgement instruments'' is retracted in
Appendix~\ref{app:corrections}: the instruments were logically equivalent, so the
coefficient was guaranteed by construction \citep{cohen1960kappa}. Nor does the
washout supply an intra-rater figure by another route: the two passes were not
retained as separate columns, the second overwrote the first in the annotation
record, and no committed artifact holds pass-one and pass-two labels side by
side. The absence is not a reporting gap that a re-run of an existing script
would close. It is a property of what the record kept. No inter-rater or
intra-rater agreement statistic exists for the 299-pair gold standard, one
cannot be recovered from the two-pass design, and every clinical result in this
thesis, the pool precision, the weak-labeller comparison, the detector's AUROC,
the danger ratings, rests on that unmeasured instrument. Board certification is
a credential, not a reliability coefficient, and it is not offered as one here.

\paragraph{No rate is claimed from the correction record.} The self-audit record
is an appendix and a demonstration, not a statistic. An attempt to compute an
error rate from it was retracted: the generating script was not saved, the
denominator advanced without the numerator, and an independent hand-code
intersected the original list on a minority of the union. The record supports
named instances, including four cases of a correcting analysis committing the
error it was correcting.

\section{Threats to Validity}
\label{sec:disc-threats}

Internal threats attack whether the measured relationships hold in this sample,
external threats whether they hold anywhere else, construct threats whether the
quantities measured are the quantities named. They are separated because they
call for different responses.

\subsection{Internal validity}

Table~\ref{tab:disc-internal} lists them and what this thesis does with each.

\begin{table}[H]
\centering
\caption[Internal threats and their status]{Internal threats and what this thesis does with each. Two are bounded by
measurement; the rest are stated rather than argued away.}
\label{tab:disc-internal}
\small
\begin{tabular}{p{0.28\linewidth}p{0.22\linewidth}p{0.36\linewidth}}
\toprule
\textbf{Threat} & \textbf{Severity} & \textbf{Status} \\
\midrule
Twelve positives; every interval wide & Severe & Stated; MDE $0.7355$ reported with every null; no significance claim on any ablation \\
\addlinespace
Training contamination of the evaluation set & Severe if unaddressed & Measured and corrected; the decontaminated ensemble is $0.5145$ against $0.4595$ contaminated, so the correction does not favour the conclusion \\
\addlinespace
Exemplar defect on action-conflict candidates & Severe for pool precision, not decisive for the central result & Checked on the defect-free stratum by a separate control model, not a rescoring of the decontaminated ACCORD: $n = 134$, $6$ positives, three-seed ensemble AUROC $0.5469$, $95\%$ bootstrap interval $0.2418$ to $0.8281$, MDE $0.8058$. ACCORD's contaminated released run reads $0.5078$ on the same rows; no decontaminated ACCORD score exists there \\
\addlinespace
Single annotator, no inter-rater or intra-rater reliability & Severe; bounds every gold-label conclusion & Stated; the earlier circular coefficient is retracted, and the two-pass design cannot recover an intra-rater figure because the passes were not retained separately \\
\addlinespace
Weak-label leakage into model inputs & Moderate to severe & Demonstrated on the rule-type head, whose target is one linear layer from an input feature that equals it; the affected metric is vacated, which leaves the head untested rather than falsified \\
\addlinespace
Extraction never validated against human annotation & Moderate to severe; no downstream error rate decomposes & Stated; first item of any continuation \\
\bottomrule
\end{tabular}
\end{table}

The contamination row deserves a sentence, because the direction of a correction
matters. A correction that improves the result it corrects invites the suspicion
that it was selected; every correction applied to a headline number here moved it
toward the negative conclusion, which is weak evidence but evidence in the right
direction.

\subsection{External validity}

Generalisation is bounded on four axes. One centre and one documentation culture.
One primary billing-code phenotype, with two weaker transfer cohorts that move
across index diagnosis rather than across sites. One era, predating current
documentation systems and any use of generative assistance in note writing. The
fourth axis is the task, and there the bound is measured rather than assumed: the
over-calling magnitude does not generalise, because a control on public text
returned a much smaller factor for the same labellers, and the part that
transfers is the mechanism, a labeller reporting a prior when the base rate is
low; the public-text maximum over-calling ratio, measured on SNLI across three
models, two prompts and four prevalence levels, is $3.5000$, well
below the $15.8196$ to $23.5959$ range measured on the clinical pool.
The positive control is external evidence of a
limited kind: it shows the instrument behaves as expected on a public benchmark,
and says nothing about whether the clinical result would reproduce elsewhere.

\subsection{Construct validity}

Four constructs are named more confidently than they are measured, which is the
gap between a construct and its operationalisation that
\citet{jacobs2021measurement} argue is where invalid conclusions originate.
``Dangerous medication contradiction'' is operationalised by four rules and one
pharmacist's judgement; a different reading of what counts as a conflict yields a
different pool and plausibly a different base rate, and the construct is not
independently validated. The detector's target is the weak label, not the
clinical event, so any statement that it ``fails to detect contradictions'' is a
statement about what it learned to predict. Ranking metrics at a $4\%$ base rate
also measure less than they appear to: AUROC is insensitive to prevalence in a
way that flatters this setting, while average precision, which is not, sits near
the base rate throughout \citep{saito2015prcurve}, and agreement statistics here
are subject to the documented paradoxes
\citep{feinstein1990paradoxes,cicchetti1990resolving,byrt1993pabak}. Matthews
correlation is used as the better-behaved default
\citep{chicco2020mcc,delgado2019kappa} rather than as a solution, and applying
that critique is not contributing it; it is thirty-six years old.

The part of that critique this evaluation set makes concrete is the exchange
rate. Each summary statistic implies a rate at which a forgone true positive is
worth some number of gained true negatives, and the rate is fixed by the class
counts rather than by anything clinical. The prevalence-adjusted bias-adjusted kappa
\citep{byrt1993pabak} has the shallowest exchange of the three most commonly
reported here, which makes it the most majority-dominated of them despite the
name, and none of the three was chosen by asking what a missed contradiction
costs relative to a false alarm. That question is clinical and it was never
posed. One worked example on this evaluation set makes the point without a
figure. The local \texttt{granite3.3:8b} few-shot labeller of
Section~\ref{sec:m-control-rq2} carries a confusion matrix of $10$ true
positives and $75$ true negatives against ACCORD's $11$ true positives and
$18$ true negatives on the same $253$ rows: it trades one additional forgone
true positive for far more gained true negatives, which is worse on the
first axis and far better on the second, and Matthews correlation prices that
trade as a net gain rather than a net loss, which is why the labeller's
coefficient of $0.0668$ is positive on this set where ACCORD's $-0.0070$ is
not, even though the class counts are what move the statistic here, not a
clinical judgement of which errors matter more. The same two confusion
matrices could support opposite readings under a statistic with a different
exchange rate, so the choice of summary statistic is not a neutral reporting
convention but a decision about which errors are worth trading, taken before
any result is seen.

Finally, the
survey's coded construct admits two defensible readings, and the gap between the
strict and permissive counts is the size of that ambiguity.

Three results would change these conclusions, and they are tests rather than
hedges. Re-annotation against corrected exemplars would convert the
pool-precision bound into an estimate. A cohort with substantially more than $12$
positives would move the minimum detectable effect down far enough that a modest
true signal, if one exists, would be visible; at present a true AUROC around
$0.65$ is indistinguishable from chance in this design. An explanation of the
selection-channel anomaly that survives a parameter sweep, unlike the five that
did not, would convert Section~\ref{sec:disc-anomaly} from a report into a
finding. None is available from the data collected here.
